\section{Related Work}
\label{back}

% \textbf{Paragraph Goal: Define the broad domain problem and the standard baseline approach used to solve it (the "old way").}
% \textbf{Specific Domain} typically relies on \textbf{Standard Architecture} to achieve \textbf{Target Goal} \citep{cite1}. However... by \textbf{Key Limitation 1} and \textbf{Key Limitation 2}... To address this, the standard operational pipeline applies \textbf{Baseline Method}, ... to a \textbf{....} via:
% \[
% \hat{x} \;=\; \mathcal{T}(x) \;\approx\; \text{some math\_op}(x, \theta),
% \]
% where....

% \textbf{Paragraph Goal: Categorize state-of-the-art improvements and identify the specific shared limitation (the "Gap")—usually that they are static or context-agnostic.}
% Recent advances in...have moved beyond // approaches, generally falling into X categories: \textit{(i) \textbf{Category A}}, which ... \citep{cite3}; and \textit{(ii) ...  Despite their effectiveness, these methods share a critical limitation: .....\citep{cite5}.

% \textbf{Paragraph Goal: Introduce the "Insight" from a parallel field (e.g., interpretability/analysis) that proves a better approaches.}
% In parallel, research in \textbf{Domain} has uncovered that \textbf{the phenomena}  \citep{cite6}. Specifically, evidence shows that \textbf{Open the phenomena} \citep{cite7}. By measuring \textbf{how we want to use this phenomena in the domain method}, one can effectively  \textbf{Improvement } from noise. This suggests that \textbf{the domain task} should not be  \textbf{break down something the induce into the phenomena }.

% \textbf{Paragraph Goal: Connect the two previous points. Explain how we operationalize the insight into a concrete metric (e.g., Sensitivity/Entropy) to bridge the gap.}
% To leverage this insight, we require a  \textbf{Component  and Importance}. We adopt a \textbf{ Framework} that evaluates the \textbf{Relevance} of  \textbf{something by leveraging the phenomena as follow}. Formally, we define ..:
% \[
% \mathcal{R}_l \;=\; f(\text{signal\_source}_l) \;\propto\; \text{Impact on Task},
% \]
% where .... This metric allows us to ...\citep{cite8}.

% \textbf{Paragraph Goal: Formulate the core problem of the paper: using the new metric to solve the  problem better than the baselines.}
% We propose a \textbf{Our method name } that ... . Unlike \textbf{Baseline Methods} that ..., we formulate the \textbf{ Problem} as:
% \[
% \min_{\mathbf{b}} \; \sum \text{some math}(b_l) \quad \text{s.t.} \quad \text{}(b_l) \ge \mathcal{R}_l,
% \]
% where.... This approach ensures that...., thereby bridging the gap between ...